% =====================================================================
\section{Part III, AI Exercise 3: Literature Synthesis and Historical Context}
% =====================================================================

\subsection{3(a)}

\begin{table}[H]
\centering
\caption{These are the values reported in Part I(b). Estimates and standard
errors are in percent per month; the table's $t$-statistics test against zero.
The reported mean market excess returns are $0.6116\%$ and $0.6952\%$ per
month, respectively. For the main sample, the report gives $t = -2.12$ for
testing the slope against the sample market premium.}
\small
\begin{tabular}{llrrrr}
\toprule
Sample & Parameter & Estimate & Standard error & $t$ & $p$ \\
\midrule
July 1969--June 2026 & $\gamma_0$ & 0.6119 & 0.2191 & 2.79 & 0.005 \\
July 1969--June 2026 & $\gamma_M$ & 0.0352 & 0.2725 & 0.13 & 0.897 \\
July 1926--June 2026 & $\gamma_0$ & 0.5160 & 0.1714 & 3.01 & 0.003 \\
July 1926--June 2026 & $\gamma_M$ & 0.2490 & 0.2315 & 1.08 & 0.282 \\
\bottomrule
\end{tabular}
\label{tab:ex3a}
\end{table}

My Part I(b) results weaken the favorable evidence of Fama and MacBeth (1973)
and contradict the Sharpe--Lintner CAPM restrictions in my sample. For July
1969--June 2026, $\hat\gamma_0 = 0.6119\%$ per month (SE $= 0.2191$,
$t = 2.79$, $p = 0.005$), significantly above the zero intercept predicted for
excess returns. Meanwhile, $\hat\gamma_M = 0.0352\%$ per month
(SE $= 0.2725$, $t = 0.13$, $p = 0.897$), providing no statistically
significant evidence of a positive beta premium. The slope is also far below
the CAPM benchmark of $0.6116\%$, with the reported benchmark test rejecting
equality at the $5\%$ level ($t = -2.12$). Thus, the estimated security market
line is nearly flat, failing to replicate their positive beta--return
relationship. The longer July 1926--June 2026 sample reinforces this pattern:
the intercept remains significantly positive, while the beta premium remains
insignificant. A positive excess-return intercept alone does not reject
Black's zero-beta CAPM, but the weak beta--return relationship undermines the
earlier favorable interpretation.

\subsection{3(b)}

The original LLM response correctly identifies Banz's (1981) size effect: small
firms earned higher average returns than their market betas could explain. It
also accurately summarizes Fama and French (1992), emphasizing the explanatory
power of size and book-to-market and the weak relation between beta and average
returns once beta variation is separated from size. However, it omits
Rosenberg, Reid, and Lanstein (1985), an important earlier study documenting
the book-to-market effect that the assignment explicitly asks about. Although
its discussion of Fama and French (1993) adds useful context, it does not fill
this omission. Overall, the response captures the main empirical challenges to
the CAPM but gives an incomplete account of the development of the value
evidence.

\subsection{3(c)}

The original LLM response does not acknowledge data snooping: size and
book-to-market may appear important partly because researchers identified and
tested these patterns using the same historical data. This concern also applies
to Part I(f), since our sample overlaps with the data used to discover these
effects, so it is not a fully independent test. However, we use previously
documented characteristics rather than search across many predictors ourselves.
Our results provide only weak evidence: size is insignificant ($t = 0.86$),
while book-to-market is marginally significant ($t = 1.95$). Lagging these
characteristics addresses their timing but does not eliminate data snooping.
Testing on a separate period after the original studies would provide stronger
independent evidence.

\subsection{3(d)}

\paragraph{LLM response to record.} International evidence shows that size and
value effects occur outside the United States, although their strength varies
across markets and samples.

\begin{tight}
\item \textbf{Direct Fama--MacBeth evidence:} Hu, Chen, Shao, and Wang (2019)
find a significant size-factor premium in Chinese stocks, while market and
value-factor premia are insignificant. Thus, international Fama--MacBeth tests
do not uniformly support both size and value. [1]
\item \textbf{Broader international evidence:} Rouwenhorst (1999) finds that
small stocks outperform large stocks and value stocks outperform growth stocks
in emerging markets. Fama and French (1998) find value outperformance in twelve
of thirteen major markets. These studies provide complementary evidence from
portfolio comparisons and factor tests; their findings should not all be
described as Fama--MacBeth regression results. [2], [3]
\end{tight}

My interpretation is that international replication reduces the likelihood that
these effects are merely artifacts of searching U.S.\ data. It is consistent
with compensation for omitted risks, but does not establish that explanation:
similar investor biases and limits to arbitrage could also produce recurring
patterns across countries.

\paragraph{Evaluation for submission.} The international evidence modestly
strengthens the plausibility of a risk-based interpretation by showing that size
and value patterns are not confined to U.S.\ data. However, it strengthens the
evidence that these patterns exist more clearly than it distinguishes risk from
mispricing, since investor biases and limits to arbitrage may also operate
internationally. The LLM appropriately acknowledges this distinction and the
variation across markets. In our Part I(f), size is insignificant ($t = 0.86$)
and book-to-market is only marginally significant ($t = 1.95$), so
international findings cannot turn our weak evidence into proof of priced
risks. They provide supporting context for the positive book-to-market
coefficient, but establishing a risk explanation would require evidence linking
returns to exposure to economically meaningful risks.

\paragraph{Sources cited in 3(d).}
\begin{sloppypar}\raggedright
[1] Hu, Chen, Shao, and Wang (2019). ``Fama--French in China: Size and Value
Factors in Chinese Stock Returns.'' \texttt{doi:10.1111/irfi.12177}\\
{[2]} Rouwenhorst (1999). ``Local Return Factors and Turnover in Emerging Stock
Markets.'' \texttt{doi:10.1111/0022-1082.00151}\\
{[3]} Fama and French (1998). ``Value versus Growth: The International
Evidence.'' \texttt{doi:10.1111/0022-1082.00080}
\end{sloppypar}
